\documentclass[11pt,a4paper]{article}
\usepackage[utf8]{inputenc}
\usepackage[T1]{fontenc}
\usepackage{amsmath,amssymb,amsthm}
\usepackage{mathtools}
\usepackage{hyperref}
\usepackage{geometry}
\usepackage{xcolor}
\usepackage{booktabs}
\geometry{margin=1in}

\hypersetup{colorlinks=true, linkcolor=blue!50!black, urlcolor=blue!50!black,
  citecolor=blue!50!black}

%% =====================================================================
%% v6 (2026-06-20).  Focus: rigorous, human-readable proofs of the three
%% congruence results that are also formally verified in Lean 4 (axiom-clean,
%% sorry-free).  Names de-eponymized; OEIS A397213 marked as a pending draft.
%% Modesty + explicit invitation for community review throughout.
%% =====================================================================

\newtheorem{theorem}{Theorem}[section]
\newtheorem{lemma}[theorem]{Lemma}
\newtheorem{proposition}[theorem]{Proposition}
\newtheorem{corollary}[theorem]{Corollary}
\theoremstyle{definition}
\newtheorem{conjecture}[theorem]{Conjecture}
\newtheorem{definition}[theorem]{Definition}
\theoremstyle{remark}
\newtheorem{remark}[theorem]{Remark}

\newcommand{\Sx}{S_{20}}
\newcommand{\modp}[1]{\ (\mathrm{mod}\ #1)}
\DeclareMathOperator{\diag}{Diag}

\title{A Weight-5 Ap\'ery-like Sequence $S_{20}(n)$:\\
  Formally Verified Supercongruences and a Conjectural Calabi--Yau Period\\
  \large (preprint, v6 --- under community review)}
\author{Xavier Callens\\
  SocrateAI Laboratory\\
  \url{https://github.com/xaviercallens/Mirror-Map-Sieve}}
\date{\today}

\begin{document}
\maketitle

\begin{abstract}
We study the integer sequence
\[
  \Sx(n) \;=\; \sum_{k=0}^{n} \binom{n}{k}^{4}\binom{n+k}{k}
  \;=\; 1,\,3,\,55,\,1155,\,29751,\,852753,\dots,
\]
a weight-$5$ (more precisely weight-$(4,1)$) Ap\'ery-like binomial sum that we
have been unable to locate in the literature or in the OEIS.  A draft OEIS
submission (proposed identifier
\href{https://oeis.org/draft/A397213}{A397213}) is \emph{pending editorial
review and has not been approved}; we therefore make no novelty claim beyond
``we did not find it.''  The purpose of this note is narrow and, we hope,
honest: we give complete, human-readable proofs of three arithmetic congruences
satisfied by $\Sx$, and we report that \emph{each of these three results has
also been verified in the Lean~4 proof assistant with no \texttt{sorry}, no
\texttt{admit}, and no axioms beyond Lean's three standard foundations}
(\texttt{propext}, \texttt{Classical.choice}, \texttt{Quot.sound}).  The three
results are:
\begin{enumerate}\itemsep0pt
  \item[(I)] the \emph{cubic supercongruence} $\Sx(p)\equiv 3 \modp{p^{3}}$ for
        every prime $p\ge 5$;
  \item[(II)] \emph{Wolstenholme's theorem} $\binom{2p}{p}\equiv 2 \modp{p^{3}}$
        for $p\ge 5$, reproved from first principles (it is the one external
        input of (I));
  \item[(III)] the \emph{Ap\'ery-style boundary congruence}
        $\Sx(p-1)\equiv 1 \modp{p}$ for every prime $p$.
\end{enumerate}
We are explicit about what is \emph{not} settled: the order-$4$ minimal linear
recurrence is computed but not formally certified for general $n$; the mod-$p^3$
refinement of (III) is an open conjecture; and the interpretation of $\Sx$ as a
Calabi--Yau period is a structural analogy, not a theorem.  We invite
correction, refutation, and independent verification.
\end{abstract}

\tableofcontents

%% =====================================================================
\section{Introduction and honest scope}
%% =====================================================================

Ap\'ery's irrationality proofs for $\zeta(2)$ and $\zeta(3)$ rest on integer
sequences satisfying three-term recurrences with remarkable arithmetic
properties.  A now-classical theme --- developed by Beukers, Zagier,
Almkvist--Zudilin, and others --- is that such sequences often coincide with
periods of one-parameter families of Calabi--Yau manifolds, and that their
$p$-adic behaviour (supercongruences) reflects this geometry.

This note concerns the specific binomial sum
\begin{equation}\label{eq:def}
  \Sx(n) \;=\; \sum_{k=0}^{n}\binom{n}{k}^{4}\binom{n+k}{k}.
\end{equation}
It is the $A=4,\ B=1$ member of the family
$\sum_k \binom{n}{k}^A\binom{n+k}{k}^B$; the Ap\'ery numbers for $\zeta(3)$ are
$A=B=2$.

\paragraph{What this paper claims, and what it does not.}
We have tried to separate rigour from speculation cleanly. The reader should
read the following table as the precise scope of the work.

\begin{center}
\begin{tabular}{@{}lll@{}}
\toprule
Statement & Status & Where \\
\midrule
First terms / closed form \eqref{eq:def} & verified (exact arithmetic) & \S\ref{sec:prelim} \\
$\Sx(p)\equiv 3\ (p^3)$, $p\ge5$ & \textbf{proved} (+ Lean) & \S\ref{sec:cubic} \\
$\binom{2p}{p}\equiv 2\ (p^3)$, $p\ge5$ & \textbf{proved} (+ Lean) & \S\ref{sec:wolstenholme} \\
$\Sx(p-1)\equiv 1\ (p)$, all $p$ & \textbf{proved} (+ Lean) & \S\ref{sec:apery} \\
$\Sx(p-1)\equiv 1\ (p^3)$ & open conjecture (numeric only) & \S\ref{sec:apery} \\
order-4 recurrence (general $n$) & computed, not formally certified & \S\ref{sec:open} \\
Calabi--Yau period interpretation & conjectural analogy & \S\ref{sec:open} \\
\bottomrule
\end{tabular}
\end{center}

\paragraph{On the Lean~4 formalization.}
All three congruence theorems below were transcribed into Lean~4 (Mathlib,
toolchain \texttt{v4.31.0}).  Building the development with \texttt{lake build}
succeeds with no errors; the command \texttt{\#print axioms} applied to each
final theorem returns only \texttt{[propext, Classical.choice, Quot.sound]}.  We
stress what this does and does not mean: it means the \emph{stated} theorems are
correct relative to Lean's kernel and Mathlib; it does \emph{not} certify the
recurrence for general $n$, nor any of the conjectural material.  We give the
human proofs in full here precisely so that they can be checked by a
mathematician without trusting the machine.  We would be grateful for scrutiny
of both.

%% =====================================================================
\section{Preliminaries}\label{sec:prelim}
%% =====================================================================

Throughout, $p$ denotes a prime and $\binom{n}{k}$ the usual binomial
coefficient, with $\binom{n}{k}=0$ for $k>n$.  We write the summand of
\eqref{eq:def} as
\[
  T(n,k) \;=\; \binom{n}{k}^{4}\binom{n+k}{k}, \qquad
  \Sx(n) \;=\; \sum_{k=0}^{n} T(n,k).
\]

We use two elementary facts repeatedly.

\begin{lemma}[Prime divides interior binomials]\label{lem:dvdchoose}
If $p$ is prime and $1\le k\le p-1$, then $p \mid \binom{p}{k}$.
\end{lemma}

\begin{proof}
From $k\binom{p}{k}=p\binom{p-1}{k-1}$ we get $p \mid k\binom{p}{k}$.  Since
$1\le k\le p-1$, the prime $p$ does not divide $k$, hence $p\mid\binom{p}{k}$.
(This is \texttt{Nat.Prime.dvd\_choose\_self} in Mathlib.)
\end{proof}

\begin{lemma}[Factorial $p$-divisibility]\label{lem:dvdfact}
For a prime $p$ and $m\in\mathbb N$, one has $p \mid m!$ if and only if
$p\le m$.
\end{lemma}

\begin{proof}
If $p\le m$ then $p$ is one of the factors of $m!$.  Conversely, if $p\mid m!$
then, as $p$ is prime, $p$ divides one of $1,2,\dots,m$, forcing $p\le m$.
(This is \texttt{Nat.Prime.dvd\_factorial}.)
\end{proof}

%% =====================================================================
\section{The cubic supercongruence $\Sx(p)\equiv 3\ (\mathrm{mod}\ p^3)$}
\label{sec:cubic}
%% =====================================================================

The main arithmetic result of this note is the following.

\begin{theorem}[Cubic supercongruence]\label{thm:cubic}
For every prime $p\ge 5$,
\[
  \Sx(p)\;\equiv\;3 \pmod{p^{3}}.
\]
\end{theorem}

The proof has two ingredients: an elementary \emph{collapse} of the sum to its
two boundary terms modulo $p^{3}$ (Proposition~\ref{prop:collapse}), and the
classical Wolstenholme congruence for $\binom{2p}{p}$
(Theorem~\ref{thm:wolstenholme}, proved independently in \S\ref{sec:wolstenholme}).
The collapse is the part specific to $\Sx$, and it is genuinely simple; the
arithmetic depth is entirely in Wolstenholme.

\subsection{The collapse to boundary terms}

\begin{proposition}[Collapse]\label{prop:collapse}
For every prime $p$,
\[
  \Sx(p)\;\equiv\; 1 + \binom{2p}{p}\pmod{p^{3}}.
\]
\end{proposition}

\begin{proof}
Split the sum \eqref{eq:def} at its endpoints:
\[
  \Sx(p) \;=\; T(p,0) \;+\; \sum_{k=1}^{p-1} T(p,k) \;+\; T(p,p).
\]
The boundary terms are explicit:
\[
  T(p,0)=\binom{p}{0}^{4}\binom{p}{0}=1,
  \qquad
  T(p,p)=\binom{p}{p}^{4}\binom{2p}{p}=\binom{2p}{p}.
\]
It remains to show that each interior term vanishes modulo $p^{3}$.  Fix
$1\le k\le p-1$.  By Lemma~\ref{lem:dvdchoose}, $p\mid\binom{p}{k}$, hence
$p^{4}\mid\binom{p}{k}^{4}$.  The companion factor $\binom{p+k}{k}$ is an
integer, so
\[
  p^{4}\ \Big|\ \binom{p}{k}^{4}\binom{p+k}{k} = T(p,k),
\]
and a fortiori $p^{3}\mid T(p,k)$.  Summing, $p^3 \mid \sum_{k=1}^{p-1}T(p,k)$,
which gives the claim.
\end{proof}

\begin{remark}
The factor of $\binom{p}{k}^{4}$ is crucial: a single power would give only
$p^{1}\mid\binom{p}{k}$, enough for a mod-$p$ statement but not mod-$p^{3}$.  The
fourth power provides $p^{4}$ ``for free,'' so the companion binomial
$\binom{p+k}{k}$ --- whatever its $p$-adic valuation --- cannot spoil
divisibility by $p^{3}$.  This is the only place the exponent $4$ in
\eqref{eq:def} is used in this section.
\end{remark}

\subsection{Proof of Theorem~\ref{thm:cubic}}

\begin{proof}
By Proposition~\ref{prop:collapse}, $\Sx(p)\equiv 1+\binom{2p}{p}\pmod{p^{3}}$.
By Wolstenholme's theorem (Theorem~\ref{thm:wolstenholme}), valid for $p\ge5$,
we have $\binom{2p}{p}\equiv 2\pmod{p^{3}}$.  Therefore
\[
  \Sx(p)\;\equiv\;1+2\;=\;3\pmod{p^{3}}. \qedhere
\]
\end{proof}

\begin{corollary}[Unconditional mod $p^2$]\label{cor:modsq}
For \emph{every} prime $p$, $\Sx(p)\equiv 3\pmod{p^{2}}$.
\end{corollary}

\begin{proof}
The collapse (Proposition~\ref{prop:collapse}) holds mod $p^{3}$ and hence mod
$p^{2}$.  Babbage's congruence $\binom{2p}{p}\equiv 2\pmod{p^{2}}$ holds for all
primes (Theorem~\ref{thm:babbage}), so $\Sx(p)\equiv 1+2=3\pmod{p^2}$.  No lower
bound on $p$ is needed here.
\end{proof}

%% =====================================================================
\section{Wolstenholme's theorem from first principles}\label{sec:wolstenholme}
%% =====================================================================

For completeness and because our Lean development does not assume it, we prove
the two congruences for the central binomial coefficient used above.  Nothing
here is new --- these are classical results of Babbage (1819) and Wolstenholme
(1862) --- but we give self-contained arguments matching the formalization.

The starting point is a Vandermonde identity.

\begin{lemma}[Central binomial as a sum of squares]\label{lem:vandermonde}
For all $p\ge 0$, $\displaystyle \binom{2p}{p}=\sum_{k=0}^{p}\binom{p}{k}^{2}.$
\end{lemma}

\begin{proof}
Chu--Vandermonde gives $\binom{2p}{p}=\sum_{k=0}^{p}\binom{p}{k}\binom{p}{p-k}$,
and $\binom{p}{p-k}=\binom{p}{k}$.
\end{proof}

\begin{theorem}[Babbage, mod $p^2$]\label{thm:babbage}
For every prime $p$, $\binom{2p}{p}\equiv 2\pmod{p^{2}}$.
\end{theorem}

\begin{proof}
By Lemma~\ref{lem:vandermonde},
$\binom{2p}{p}=\binom{p}{0}^2+\sum_{k=1}^{p-1}\binom{p}{k}^2+\binom{p}{p}^2
 = 2+\sum_{k=1}^{p-1}\binom{p}{k}^{2}$.
For $1\le k\le p-1$, Lemma~\ref{lem:dvdchoose} gives $p\mid\binom{p}{k}$, hence
$p^{2}\mid\binom{p}{k}^{2}$.  Thus the interior sum is divisible by $p^{2}$ and
$\binom{2p}{p}\equiv 2\pmod{p^2}$.
\end{proof}

To upgrade to mod $p^{3}$ we need the harmonic input.  Work in the field
$\mathbb F_p=\mathbb Z/p\mathbb Z$.

\begin{lemma}[Pascal residues]\label{lem:pascal}
For a prime $p$ and $0\le j\le p-1$, $\binom{p-1}{j}\equiv(-1)^{j}\pmod p$.
\end{lemma}

\begin{proof}
Induction on $j$.  For $j=0$, $\binom{p-1}{0}=1$.  For the step, Pascal's rule
gives $\binom{p}{j+1}=\binom{p-1}{j}+\binom{p-1}{j+1}$.  By
Lemma~\ref{lem:dvdchoose}, $p\mid\binom{p}{j+1}$ when $1\le j+1\le p-1$, so in
$\mathbb F_p$ we have $\binom{p-1}{j+1}\equiv-\binom{p-1}{j}\equiv
-(-1)^{j}=(-1)^{j+1}$.
\end{proof}

\begin{lemma}[Harmonic core]\label{lem:harmonic}
For a prime $p\ge 5$, $\displaystyle\sum_{k=1}^{p-1}\frac{1}{k^{2}}\equiv 0
\pmod p$, where $1/k$ denotes the inverse of $k$ in $\mathbb F_p$.
\end{lemma}

\begin{proof}
As $k$ ranges over $1,\dots,p-1$, the inverses $k^{-1}$ range over the same set
$\mathbb F_p^{\times}$, so $\sum_{k=1}^{p-1}k^{-2}=\sum_{x\in\mathbb
F_p^\times}x^{2}=\sum_{x\in\mathbb F_p}x^{2}$ (the $x=0$ term is $0$).  For any
finite field $\mathbb F_q$ and exponent $i$ with $0<i<q-1$, one has
$\sum_{x\in\mathbb F_q}x^{i}=0$ (a standard fact: the sum of all $i$-th powers
over a finite field vanishes unless $(q-1)\mid i$).  Here $q=p$, $i=2$, and
$2<p-1$ exactly when $p\ge5$, so the sum is $0$.
\end{proof}

\begin{theorem}[Wolstenholme, mod $p^3$]\label{thm:wolstenholme}
For every prime $p\ge 5$, $\binom{2p}{p}\equiv 2\pmod{p^{3}}$.
\end{theorem}

\begin{proof}
By Lemma~\ref{lem:vandermonde} and the boundary evaluation as in
Theorem~\ref{thm:babbage},
\[
  \binom{2p}{p}-2 \;=\; \sum_{k=1}^{p-1}\binom{p}{k}^{2}.
\]
For $1\le k\le p-1$ write $\binom{p}{k}=\tfrac{p}{k}\binom{p-1}{k-1}$, which is
the integer identity $k\binom{p}{k}=p\binom{p-1}{k-1}$ solved for
$\binom{p}{k}$; set $d_k:=\binom{p}{k}/p=\tfrac1k\binom{p-1}{k-1}\in\mathbb Z$
(an integer by Lemma~\ref{lem:dvdchoose}).  Then $\binom{p}{k}^{2}=p^{2}d_k^{2}$
and
\[
  \binom{2p}{p}-2 \;=\; p^{2}\sum_{k=1}^{p-1} d_k^{2}.
\]
Thus $\binom{2p}{p}\equiv 2\pmod{p^{3}}$ if and only if
$p\mid\sum_{k=1}^{p-1}d_k^{2}$.  Reduce $d_k$ modulo $p$.  From
$k\,d_k=\binom{p-1}{k-1}$ and Lemma~\ref{lem:pascal},
$k\,d_k\equiv(-1)^{k-1}\pmod p$, so in $\mathbb F_p$
\[
  d_k\equiv (-1)^{k-1}k^{-1},\qquad
  d_k^{2}\equiv k^{-2}.
\]
Hence $\sum_{k=1}^{p-1}d_k^{2}\equiv\sum_{k=1}^{p-1}k^{-2}\equiv 0\pmod p$ by
Lemma~\ref{lem:harmonic} (using $p\ge5$).  Therefore
$p\mid\sum_k d_k^2$ and $\binom{2p}{p}\equiv 2\pmod{p^3}$.
\end{proof}

\begin{remark}
The hypothesis $p\ge 5$ enters in exactly one place: Lemma~\ref{lem:harmonic}
needs $2<p-1$.  For $p=3$ one checks directly $\binom{6}{3}=20\equiv 2\pmod 9$
but $20\equiv 20\not\equiv2\pmod{27}$, so the mod-$p^3$ statement genuinely
fails at $p=3$; the mod-$p^2$ statement (Theorem~\ref{thm:babbage}) survives.
\end{remark}

%% =====================================================================
\section{The Ap\'ery-style boundary congruence}\label{sec:apery}
%% =====================================================================

The second sequence-specific result concerns the value at $p-1$.

\begin{theorem}[Ap\'ery-style congruence mod $p$]\label{thm:apery}
For every prime $p$, $\;\Sx(p-1)\equiv 1\pmod p$.
\end{theorem}

The mechanism mirrors the collapse of \S\ref{sec:cubic}, but now the
divisibility comes from the \emph{companion} binomial rather than from
$\binom{p-1}{k}$.

\begin{lemma}[Companion divisibility]\label{lem:companion}
For a prime $p$ and $1\le k\le p-1$, $\;p\mid\binom{p-1+k}{k}$.
\end{lemma}

\begin{proof}
The factorial identity for binomial coefficients gives
\[
  \binom{p-1+k}{k}\cdot k!\cdot (p-1)! \;=\; (p-1+k)!,
\]
since $(p-1+k)-k=p-1$.  As $1\le k\le p-1$ we have $p\le p-1+k$, so by
Lemma~\ref{lem:dvdfact} the prime $p$ divides the right-hand side $(p-1+k)!$.
On the other hand $k<p$ and $p-1<p$, so by Lemma~\ref{lem:dvdfact} again $p\nmid
k!$ and $p\nmid(p-1)!$.  Since $p$ is prime and divides the product
$\binom{p-1+k}{k}\,k!\,(p-1)!$ but neither $k!$ nor $(p-1)!$, it must divide the
remaining factor $\binom{p-1+k}{k}$.
\end{proof}

\begin{proof}[Proof of Theorem~\ref{thm:apery}]
Split the sum at $k=0$:
\[
  \Sx(p-1) \;=\; T(p-1,0) \;+\; \sum_{k=1}^{p-1} T(p-1,k),
  \qquad T(p-1,0)=\binom{p-1}{0}^4\binom{p-1}{0}=1.
\]
For each interior index $1\le k\le p-1$,
$T(p-1,k)=\binom{p-1}{k}^{4}\binom{p-1+k}{k}$, and
Lemma~\ref{lem:companion} gives $p\mid\binom{p-1+k}{k}$, hence $p\mid
T(p-1,k)$.  Therefore $\sum_{k=1}^{p-1}T(p-1,k)\equiv 0\pmod p$ and
$\Sx(p-1)\equiv 1\pmod p$.
\end{proof}

\begin{conjecture}[Strong Ap\'ery-style congruence]\label{conj:apery3}
For every prime $p\ge 5$, $\;\Sx(p-1)\equiv 1\pmod{p^{3}}$.
\end{conjecture}

We have verified Conjecture~\ref{conj:apery3} for all primes $5\le p\le 200$
(see the reproducibility report).  We do \emph{not} have a proof; the mod-$p$
argument above does not extend naively, because $\binom{p-1}{k}^4\equiv 1\pmod
p$ controls only the leading $p$-adic order.  A proof would presumably require a
harmonic-sum analysis in the spirit of \S\ref{sec:wolstenholme} (likely
involving $\sum 1/k$ and $\sum 1/k^2$ to higher order), and we would welcome one.

%% =====================================================================
\section{Numerical and formal verification}\label{sec:verification}
%% =====================================================================

All numerical claims are reproduced by a single self-contained script,
\texttt{python/verify\_all.py} (standard library only), whose checks include:
the first $18$ terms against an independent transcription; the cubic
supercongruence and Wolstenholme/Babbage congruences for all primes up to
$200$; the Ap\'ery-style congruence mod $p$ (and the conjectural mod $p^3$) up
to $200$; the collapse identity of Proposition~\ref{prop:collapse}; and the
diagonal evaluation discussed in \S\ref{sec:open}.  The script terminates with
exit code $0$ and the message \emph{``ALL NUMERICAL CHECKS PASSED.''}

Theorems~\ref{thm:cubic}, \ref{thm:wolstenholme} (with
Theorem~\ref{thm:babbage}), and \ref{thm:apery} have been formalized in Lean~4
(Mathlib, toolchain \texttt{v4.31.0}) as, respectively,
\texttt{Supercongruence.supercongruence\_unconditional},
\texttt{Wolstenholme.wolstenholme} / \texttt{Wolstenholme.babbage}, and
\texttt{AperyCongruence.apery\_congruence}.  Each compiles under
\texttt{lake build} and, under \texttt{\#print axioms}, depends only on
\texttt{propext}, \texttt{Classical.choice}, and \texttt{Quot.sound}.  The full
audit is recorded in \texttt{VERIFICATION\_REPORT.md}.

%% =====================================================================
\section{What remains open}\label{sec:open}
%% =====================================================================

We close by stating clearly the parts of the broader picture that are
\emph{not} theorems in this paper, in the hope that experts will engage with
them.

\paragraph{Linear recurrence.}
Numerical linear algebra (a $\mathbb Q$-nullspace computation) indicates that
$\Sx(n)$ satisfies a minimal linear recurrence of order $4$ and polynomial
degree $13$, with an order-$5$, degree-$9$ left-multiple.  We have verified
small base cases in Lean, but the general-$n$ statement currently rests on a
creative-telescoping (Zeilberger) certificate that we have \emph{not}
re-derived inside a formally verified setting.  A verified certificate would be
a valuable contribution.

\paragraph{Diagonal representation.}
We observe the identity, verified for $0\le n\le 6$ and provable by elementary
power-series extraction,
\[
  \Sx(n) \;=\; \diag\!\left[\frac{1}{(1-x_1)(1-x_2)(1-x_3)(1-x_4)(1-x_5)
  - x_1x_2x_3x_4}\right],
\]
i.e.\ $\Sx(n)$ is the coefficient of $(x_1\cdots x_5)^n$.  Indeed expanding
$1/(D-M)=\sum_{r\ge0}M^r/D^{r+1}$ with $D=\prod_i(1-x_i)$, $M=x_1x_2x_3x_4$ and
extracting the diagonal yields exactly $\sum_{r}\binom{n}{r}^4\binom{n+r}{r}$.
The \emph{symmetric} five-variable candidate does not represent $\Sx$; its
diagonal is $2^n$.

\paragraph{Calabi--Yau period.}
The order of the (conjectural) minimal Picard--Fuchs operator is $4$, which
would be consistent with a one-parameter family whose period geometry is that of
a Calabi--Yau $3$-fold.  We computed mirror-map coefficients $q_d$ via exact
rational arithmetic and found them to be integers for $1\le d\le 16$,
consistent with Lian--Yau integrality.  We emphasize that this is
\emph{evidence and analogy}: we do not prove that $\Sx$ is the period of a
specific Calabi--Yau manifold, and the precise geometry (including its
dimension) deserves expert scrutiny.

\paragraph{Invitation.}
This is a preprint released specifically to invite review.  Corrections,
counterexamples, a proof of Conjecture~\ref{conj:apery3}, a verified recurrence
certificate, an identification (or refutation) of the Calabi--Yau
interpretation, or a prior appearance of $\Sx$ in the literature would all be
very welcome.  Please open an issue or pull request at the repository.

%% =====================================================================
\section*{Acknowledgements}
%% =====================================================================
We thank the OEIS editors and the Lean/Mathlib community.  Any errors are the
author's own, and the work is offered in a spirit of provisional, correctable
inquiry rather than finality.

\begin{thebibliography}{9}
\bibitem{apery} R.~Ap\'ery, \emph{Irrationalit\'e de $\zeta(2)$ et
  $\zeta(3)$}, Ast\'erisque \textbf{61} (1979), 11--13.
\bibitem{babbage} C.~Babbage, \emph{Demonstration of a theorem relating to
  prime numbers}, Edinburgh Phil.\ J.\ \textbf{1} (1819), 46--49.
\bibitem{wolstenholme} J.~Wolstenholme, \emph{On certain properties of prime
  numbers}, Quart.\ J.\ Pure Appl.\ Math.\ \textbf{5} (1862), 35--39.
\bibitem{lianyau} B.~H.~Lian and S.-T.~Yau, \emph{Arithmetic properties of
  mirror map and quantum coupling}, Comm.\ Math.\ Phys.\ \textbf{176} (1996),
  163--191.
\bibitem{almkvistzudilin} G.~Almkvist, C.~van Enckevort, D.~van Straten,
  W.~Zudilin, \emph{Tables of Calabi--Yau equations}, arXiv:math/0507430.
\bibitem{mathlib} The mathlib Community, \emph{The Lean mathematical library},
  CPP 2020.
\bibitem{repo} X.~Callens, \emph{Mirror Map Sieve} (software and Lean
  formalization), \url{https://github.com/xaviercallens/Mirror-Map-Sieve},
  DOI:\,10.5281/zenodo.20747943, 2026.
\end{thebibliography}

\end{document}
